% =====================================================================
% 2. Learning and evaluation framework (shared protocol; T7 by MM, T8 shared)
% =====================================================================
\section{Learning and evaluation framework}\label{sec:framework}

Every classifier of the project shares the data, the loss, the optimiser, the metric and the
comparison rules, so that differences between models come from the architecture. The evaluation
module was written by Mar\'ia (task T7) and the training loop by both of us (T8).

\subsection{Output, loss and optimisation}

A model maps a record $x \in \mathbb{R}^{12\times1000}$ to five real scores (logits) $z_1,\dots,z_5$,
one per superclass. Each class is an independent yes/no question, so its probability is a sigmoid,
$p_c = \sigma(z_c) = 1/(1+e^{-z_c})$, and the loss is the binary cross-entropy averaged over the five
classes and the records of a batch:
\begin{equation}
\mathcal{L} = -\frac{1}{5}\sum_{c=1}^{5}\Big[\,y_c \log \sigma(z_c) + (1-y_c)\log\big(1-\sigma(z_c)\big)\Big],
\qquad y_c \in \{0,1\}.
\end{equation}
It is the negative log-likelihood of five Bernoulli variables. The sigmoid is not part of the network:
PyTorch's \texttt{BCEWithLogitsLoss} applies it inside the loss in a numerically stable way, and the
probabilities are computed only when evaluating. A softmax would be wrong here, because it forces the
five probabilities to sum to one and makes the classes compete, while
\edaFracMultilabelAllPct{} of the records have several of them.

The weights are trained with Adam, whose hyper-parameters ``typically require little tuning''
\cite[p.~1]{kingma2015adam}, with the same defaults for every family (learning rate $10^{-3}$,
batches of 64 records, at most 50 epochs). Each epoch the model is evaluated on the validation fold;
training stops when the validation macro AUROC has not improved for five epochs (\emph{early
stopping}) and the weights of the best epoch are restored. Early stopping is a regulariser: it limits
how far the weights move towards fitting the training set once validation stops improving.

\subsection{Metric and thresholds}

The area under the ROC curve (AUROC) of one class is the probability that a randomly chosen positive
record receives a higher score than a randomly chosen negative one; 0.5 is chance and 1 a perfect
ranking. It needs no threshold and does not change with the prevalence of the class, which makes it
suitable for imbalanced classes. As in the benchmark, the main metric is the macro AUROC, ``obtained
by averaging class-wise AUCs over all classes'' \cite[p.~1522]{strodthoff2021deep}, so that HYP
weighs as much as NORM. A decision rule is still needed for the F1 score: for each class, the
threshold that maximises its F1 on the validation fold is applied unchanged to the test fold.

\subsection{Uncertainty and comparisons}

A single test number hides two sources of variation: the random initialisation and order of the
batches (\emph{seed}) and the finite test set. Every configuration is therefore trained with three
seeds (42, 43, 44), and the reports give the mean and standard deviation over seeds. The finite test
set is handled with the bootstrap: the test patients are resampled with replacement 1000 times
(all records of a drawn patient enter together, because records of the same patient are correlated),
the metric is recomputed on each resample, and the 2.5 and 97.5 percentiles form a 95\,\% interval.
Resampling records instead of patients would give intervals that are too narrow.

To compare two models we use the \emph{paired} bootstrap: both are evaluated on the same resamples,
and the interval of the \emph{difference} is computed. Two models with overlapping individual intervals
can still differ consistently, because most of the variation of a resample affects both in the same
way. For one seed, a model is called better only when the interval of the difference lies wholly above
zero; the tables report in how many of the three seeds this happens (``seeds better / worse''). An effect
is called consistent only when it holds in all three seeds; one or two seeds of three are reported as weak
evidence. The bootstrap interval reflects the finite test set, not the training randomness; the seeds
cover the latter.

\subsection{Model selection}

All ablations and every choice of a final model use the validation fold only. The selection rule of
each block (the configuration with the highest mean validation macro AUROC) was written down before
its runs, and the test fold was evaluated once per final model; the training code refuses a second
test run of the same model and seed. A consequence, discussed in each block, is that a rule fixed in
advance sometimes selects a configuration that is statistically indistinguishable from its baseline.
We keep the rule's choice, because changing it after seeing the results would be selection on noise,
and we say so.
